\documentclass[12pt]{article}

\input{../../_assets/preambulo.tex}

\newcommand{\norma}[1]{\lVert #1 \rVert}
\newcommand{\normagenerica}{\lVert \cdot \rVert}
\newcommand{\dual}[1]{#1^{*}}
\newcommand{\prodescalar}[2]{\langle #1, #2 \rangle}

\begin{document}

    % 1. Foto de fondo
    % 2. Título
    % 3. Encabezado Izquierdo
    % 4. Color de fondo
    % 5. Coord x del titulo
    % 6. Coord y del titulo
    % 7. Fecha

    \input{../../_assets/portada}
    \portadaExamen{ffccA4.jpg}{Análisis Funcional\\Examen XV}{Análisis Funcional. Examen XV}{MidnightBlue}{-8}{28}{2026}{Daniel Morán Sánchez}

    \begin{description}
        \item[Asignatura] Análisis Funcional.
        \item[Curso Académico] 2024-25.
        \item[Grado] Doble Grado en Ingeniería Informática y Matemáticas.
        \item[Grupo] Único.
        \item[Profesor] David Arcoya Álvarez.
        \item[Descripción] Examen Extraordinario.
        % \item[Fecha] 19 de Noviembre de 2025.
        % \item[Duración] 2 horas.
    
    \end{description}
    \newpage


    % ------------------------------------
    
    \begin{ejercicio}[2 puntos]
Enunciado y demostración del \textbf{teorema de la aplicación abierta}.
    \end{ejercicio}
    
    \begin{ejercicio}[5 puntos]~
    \begin{enumerate}
    \item \textbf{[1 punto]}  
    Sea $E$ un espacio vectorial y $B : E \times E \to \mathbb{R}$ una aplicación bilineal,
    simétrica
    \[
    B(x,y)=B(y,x), \qquad \forall x,y\in E,
    \]
    y semidefinida positiva
    \[
    B(x,x)\ge 0, \qquad \forall x\in E.
    \]
    Prueba la desigualdad de Cauchy--Schwarz:
    \[
    |B(x,y)| \le B(x,x)^{1/2} B(y,y)^{1/2},
    \qquad \forall x,y\in E.
    \]
    
    \item \textbf{[2 puntos]}  
    Sea $E$ un espacio de Banach y $T : E \to E^\ast$ una aplicación lineal verificando
    \[
    \langle Tx,x\rangle_{E^\ast,E} \ge 0,
    \qquad \forall x\in E.
    \]
    Prueba que $T$ es acotado.
    
    \item \textbf{[2 puntos]}  
    Sea $H$ un espacio de Hilbert con su producto escalar $\langle\cdot,\cdot\rangle$
    y sea $T : H \to H$ una aplicación lineal biyectiva verificando
    \[
    \langle Tx,x\rangle \ge 0,
    \qquad
    \langle Tx,y\rangle = \langle x,Ty\rangle,
    \qquad \forall x,y\in H.
    \]
    Prueba que
    \[
    [x,y] := \langle Tx,y\rangle,
    \qquad \forall x,y\in H,
    \]
    define un producto escalar equivalente en $H$.
    \end{enumerate}
    \end{ejercicio}
    
    \begin{ejercicio}[3 puntos]
    Sea $\Phi : \mathbb{R}\to\mathbb{R}$ una función continua, no constantemente cero, con soporte compacto contenido en $(0,1)$  y con
    \[
    \Phi(x)\ge 0 \quad \forall x\in\mathbb{R},
    \]
    
    Para cada $n\in\mathbb{N}$ definimos
    \[
    f_n(x) := \Phi(x+n), \qquad \forall x\in\mathbb{R}.
    \]
    
    \begin{enumerate}
    \item \textbf{[1 punto]}  
    Prueba que la sucesión $\{f_n\}$ converge uniformemente a cero en compactos de $\mathbb{R}$
    y que para cada $p\in[1,\infty)$ la norma $\|f_n\|_{L^p(\mathbb{R})}$ es una constante
    estrictamente positiva independiente de $n\in\mathbb{N}$.
    
    \item \textbf{[2 puntos]}  
    Prueba que la sucesión $\{f_n\}$ converge a cero $\ast$--débilmente en $L^\infty(\mathbb{R})$,
    es decir, en la topología $\sigma(L^\infty(\mathbb{R}),L^1(\mathbb{R}))$.
    \end{enumerate}
    \end{ejercicio}
    
    \newpage
    \setcounter{ejercicio}{0} % Reiniciar contador de ejercicios

    \begin{ejercicio}[2 puntos]
    Enunciado y demostración del teorema de la aplicación abierta.\\

    Consultar los apuntes de teoría.
    \end{ejercicio}
    
    \begin{ejercicio}[5 puntos]~
    \begin{enumerate}
    \item \textbf{[1 punto]}  
    Sea $E$ un espacio vectorial y $B : E \times E \to \mathbb{R}$ una aplicación bilineal,
    simétrica
    \[
    B(x,y)=B(y,x), \qquad \forall x,y\in E,
    \]
    y semidefinida positiva
    \[
    B(x,x)\ge 0, \qquad \forall x\in E.
    \]
    Prueba la desigualdad de Cauchy--Schwarz:
    \[
    |B(x,y)| \le B(x,x)^{1/2} B(y,y)^{1/2},
    \qquad \forall x,y\in E.
    \]

    Sean $x,y\in E$ arbitrarios. Como $B$ es bilineal y simétrica, para todo $t\in\mathbb{R}$ se tiene
    \[
    0 \le B(ty+x,ty+x)
    = t^2 B(y,y) + 2t B(x,y) + B(x,x).
    \]
    
    El polinomio cuadrático en la variable $t$
    \[
    p(t)=t^2 B(y,y)+2tB(x,y)+B(x,x)
    \]
    es no negativo para todo $t\in\mathbb{R}$, luego su discriminante es menor o igual que cero:
    \[
    \Delta = 4B(x,y)^2 - 4B(x,x)B(y,y) \le 0.
    \]
    
    Por tanto,
    \[
    B(x,y)^2 \le B(x,x)B(y,y),
    \]
    y tomando raíces cuadradas obtenemos
    \[
    |B(x,y)| \le B(x,x)^{1/2} B(y,y)^{1/2},
    \qquad \forall x,y\in E,
    \]
    lo que concluye la demostración.

    
    \item \textbf{[2 puntos]}  
    Sea $E$ un espacio de Banach y $T : E \to E^\ast$ una aplicación lineal verificando
    \[
    \langle Tx,x\rangle_{E^\ast,E} \ge 0,
    \qquad \forall x\in E.
    \]
    Prueba que $T$ es acotado.

    Primero, usando la positividad de $T$, observamos que para todo $x,y\in E$ se tiene
    \[
    \langle T(x-y),x-y\rangle \ge 0,
    \]
    es decir,
    \[
    \langle Tx,x\rangle - \langle Tx,y\rangle - \langle Ty,x\rangle + \langle Ty,y\rangle \ge 0.
    \tag{1}
    \]
    Demostrar T acotado equivale a demostrar T continua, usaremos el teorema de la gráfica cerrada, para ello demostraremos que
    \begin{center}
        ``toda sucesión $\{x_n\}\subset E$ verificando $
            \left\{\begin{array}{l}
                \{x_n\} \stackrel{(E,\norma{\cdot}_E)}{\longrightarrow} x \\
                \{Tx_n\} \stackrel{(E^\ast,\norma{\cdot}_{E^\ast)}}{\longrightarrow} f
            \end{array}\right.$ cumple $Tx=f$''
    \end{center}

    Sea pues $\{x_n\}\stackrel{\norma{\cdot}_E}{\longrightarrow} x$ y $\{Tx_n\}\stackrel{\norma{\cdot}_{E^\ast}}{\longrightarrow} f$. Por (1), para todo $y\in E$ se tiene
    
    \[
    \langle Tx_n, x_n\rangle
    - \langle Ty, x_n\rangle
    - \langle Tx_n, y\rangle
    + \langle Ty, y\rangle
    =
    \langle Tx_n - Ty,\, x_n - y\rangle
    \ge 0.
    \]
    
    Pasando al límite cuando $n\to\infty$ (usando la continuidad del producto dual),
    obtenemos
    
    \[
    \langle f, x\rangle
    - \langle Ty, x\rangle
    - \langle f, y\rangle
    + \langle Ty, y\rangle
    \ge 0,
    \qquad \forall y\in E.
    \]
    
    Consecuentemente,
    
    \[
    \langle f - Ty,\, x - y\rangle \ge 0,
    \qquad \forall y\in E.
    \]

    
    Tomando ahora $y = x + az$, con $a\in\mathbb{R}$ y $z\in E$, se obtiene
    \[
    \langle f-T(x+az),\, x-(x+az)\rangle \ge 0,
    \]
    esto es,
    \[
    a \langle f-Tx, z\rangle \le a^2 \langle Tz,z\rangle,
    \qquad \forall z\in E,\ \forall a\in\mathbb{R}.
    \]
    
    Dividiendo por $a$ y haciendo tender $a\to 0^+$ y $a\to 0^-$, se deduce
    \[
    \begin{cases}
    \langle f-Tx,\, z\rangle \le \displaystyle\lim_{a\to 0^+} a\,\langle Tz,z\rangle =0,\\[6pt]
    \langle f-Tx,\, z\rangle \ge \displaystyle\lim_{a\to 0^-} a\,\langle Tz,z\rangle=0,
    \end{cases}
    \]
    y por tanto
    \[
    \langle f-Tx,\, z\rangle=0,
    \qquad \forall z\in E.
    \]

    Por tanto,
    \[
    f = Tx.
    \]
    
    Hemos probado que el grafo de $T$ es cerrado. Como $E$ es Banach y $E^\ast$ también lo es,
    por el \textbf{Teorema de la Gráfica Cerrada}, se concluye que $T$ es continuo,
    es decir, $T$ es acotado.
    
    \item \textbf{[2 puntos]}  
    Sea $H$ un espacio de Hilbert con su producto escalar $\langle\cdot,\cdot\rangle$
    y sea $T : H \to H$ una aplicación lineal biyectiva verificando
    \[
    \langle Tx,x\rangle \ge 0,
    \qquad
    \langle Tx,y\rangle = \langle x,Ty\rangle,
    \qquad \forall x,y\in H.
    \]
    Prueba que
    \[
    [x,y] := \langle Tx,y\rangle,
    \qquad \forall x,y\in H,
    \]
    define un producto escalar equivalente en $H$. \\
    \end{enumerate}


    \noindent Es fácil comprobar que $[\cdot,\cdot]$ es bilineal.  
    También es simétrica, pues
    \[
    [x,y] = \langle Tx,y\rangle
    = \langle x,Ty\rangle
    = \langle Ty,x\rangle
    = [y,x].
    \]
    
    \noindent Además,
    \[
    [x,x] = \langle Tx,x\rangle \ge 0.
    \]
    
    \noindent Por tanto, por el apartado a) se cumple la desigualdad de Cauchy--Schwarz:
    \[
    |[x,y]| \le [x,x]^{1/2}[y,y]^{1/2},
    \qquad \forall x,y\in H.
    \]
    
    \noindent Recordando el isomorfismo de Riesz
    \[
    H \longrightarrow H^\ast,
    \qquad x \longmapsto \langle x,\cdot\rangle,
    \]
    y que
    \[
    \langle\!\langle \langle x,\cdot\rangle ,y\rangle\!\rangle_{H^\ast,H}
    = \langle x,y\rangle,
    \qquad \forall x,y\in H,
    \]
    podemos usar el apartado b) para ver que la aplicación lineal
    \[
    H \longrightarrow H^\ast,
    \qquad x \longmapsto \langle Tx,\cdot\rangle
    \]
    es acotada, lo que equivale a que
    \[
    \norma{Tx} \le K\norma{x},
    \qquad \forall x\in H.
    \]
    
    Por el 1er corolario del teorema de la aplicación abierta, al ser $T$ biyectiva,
    se deduce también que $T^{-1}$ es continua. \\

    La comprobación de que $[x,x]=0 \Longleftrightarrow x=0$ la veremos mas adelante\\
    
    Sea $\|\cdot\|$ la norma asociada al producto escalar $\langle\cdot,\cdot\rangle$
    y definimos
    \[
    |x| := [x,x]^{1/2} \text{ (será la norma asociada al producto escalar } [\cdot,\cdot])
    \]
    
    Por la desigualdad de Cauchy--Schwarz para $\langle\cdot,\cdot\rangle$,
    \[
    |x|^2 = [x,x] = \langle Tx,x\rangle
    \le \norma{Tx}\,\norma{x}
    \overset{T \text{ acotada}}{\le} K\norma{x}^2,
    \]
    luego
    \[
    |x| \le \sqrt{K}\,\norma{x},
    \qquad \forall x\in H.
    \]
    
    Para la desigualdad recíproca, usando de nuevo Cauchy--Schwarz para $[\cdot,\cdot]$:
    \[
    \norma{x}^2
    = \langle x,x\rangle
    = [T^{-1}x,x]
    \le \|T^{-1}x\|\,|x|
    \overset{T^{-1} \text{ continua}}{\le} \sqrt{K}\,\|T^{-1}\|\,\norma{x}\,|x|.
    \]
    
    Por tanto
    \[
    \norma{x}
    \le \sqrt{K}\,\|T^{-1}\|\,|x|.
    \]
    
    Juntando ambas desigualdades:
    \[
    |x|
    \le \sqrt{K}\,\,\norma{x}
    \le \sqrt{K} \cdot (\sqrt{K} \|T^{-1}\| |x|) = K\,\|T^{-1}\|\,|x|
    \qquad \forall x\in H.
    \]
    
    De aquí se deduce:
   
    \begin{enumerate}
    \item $[x,x]=0 \Longleftrightarrow |x|=0 \Longleftrightarrow x=0,$ luego $[\cdot,\cdot]$ es un producto escalar y $|\cdot|$ es la norma asociada al producto escalar.
    \item Las normas $\|\cdot\|$ y $|\cdot|$ son equivalentes.
    \end{enumerate}
    \end{ejercicio}
    
    \begin{ejercicio}[4 puntos ]
    Sea $\Phi : \mathbb{R}\to\mathbb{R}$ una función continua, no constantemente cero, con soporte compacto contenido en $(0,1)$  y con
    \[
    \Phi(x)\ge 0 \quad \forall x\in\mathbb{R},
    \]
    
    Para cada $n\in\mathbb{N}$ definimos
    \[
    f_n(x) := \Phi(x+n), \qquad \forall x\in\mathbb{R}.
    \]
    
    \begin{enumerate}
    \item \textbf{[1 punto]}  
    Prueba que la sucesión $\{f_n\}$ converge uniformemente a cero en compactos de $\mathbb{R}$
    y que para cada $p\in[1,\infty)$ la norma $\|f_n\|_{L^p(\mathbb{R})}$ es una constante
    estrictamente positiva independiente de $n\in\mathbb{N}$.
    
    \item \textbf{[2 puntos]}  
    Prueba que la sucesión $\{f_n\}$ converge a cero $\ast$--débilmente en $L^\infty(\mathbb{R})$,
    es decir, en la topología $\sigma(L^\infty(\mathbb{R}),L^1(\mathbb{R}))$.
    \end{enumerate}


    \begin{enumerate}
    \item \textbf{Convergencia en compactos y norma $L^p$.}
        Sea $K$ un subconjunto compacto de $\mathbb{R}$. Por ser acotado, existirá $n_0\in\mathbb{N}$ tal que
    \[
    x \ge -n_0, \qquad \forall x\in K.
    \]
    Así,
    \[
    n+x \ge n-n_0 \ge 1, \qquad \forall n\ge n_0+1.
    \]
    Como el soporte de $\Phi$ está contenido en $(0,1)$, se tiene
    \[
    f_n(x)=\Phi(n+x)=0 \longrightarrow 0,
    \qquad \forall x\in K, \quad \forall n\ge n_0+1.
    \]
    Por tanto, $f_n\to 0$ uniformemente en compactos de $\mathbb{R}$.
    
    Por otra parte, usando nuevamente que el soporte de $\Phi$ está contenido en $(0,1)$,
    para todo $p\in[1,\infty)$ se tiene
    \[
    \|f_n\|_{L^p(\mathbb{R})}^p
    = \int_{-\infty}^{+\infty} |f_n(x)|^p\,dx
    \overset{\Phi\ge 0}{=} \int_{-\infty}^{+\infty} \Phi(n+x)^p\,dx.
    \overset{y=n+x}{=} \int_{-\infty}^{+\infty} \Phi(y)^p\,dy
    = \int_0^1 \Phi(y)^p\,dy,
    \]
    que es una constante independiente de $n\in\mathbb{N}$ y estrictamente positiva
    porque $\Phi$ no es idénticamente nula.
    
    \item \textbf{Convergencia $\ast$--débil.}
        \begin{description}
        \item[Opción 1.] Queremos ver que $\langle f_n, g \rangle \to 0$ para toda $g \in L^1(\mathbb{R})$.
            Sea $g \in L^1(\mathbb{R})$. Dado $\varepsilon > 0$, existe una función continua con soporte compacto $\psi \in C_c(\mathbb{R})$ tal que $\norma{g - \psi}_{L^1} < \varepsilon$ (Porque $C_c^\infty(\mathbb{R}) \subset C_c(\mathbb{R})$ es denso en $L^1(\mathbb{R})$, entonces $C_c(\mathbb{R})$ es denso en $L^1(\mathbb{R})$.). Entonces:
            \[
            |\langle f_n, g \rangle| \le |\langle f_n, g - \psi \rangle| + |\langle f_n, \psi \rangle| \le \norma{f_n}_{L^\infty} \norma{g - \psi}_{L^1} + |\langle f_n, \psi \rangle|.
            \]
            Como $\norma{f_n}_{L^\infty} = \norma{\Phi}_{L^\infty}$ está acotada y $\norma{g - \psi}_{L^1} < \varepsilon$, el primer término es pequeño.
            
            Para el segundo, sea $K$ $\psi$ el soporte de $\psi$, por el apartado (a) sabemos que $f_n \to 0$ uniformemente en $K$, luego
                
            \[
            \left| \int_{\mathbb{R}} f_n \psi| \right|
            = \left|\int_K f_n \psi \right|
            \le \sup_{x \in K} |f_n(x)| 
            \int_K |\psi(x)|\, dx 
            \xrightarrow[n \to \infty]{} 0.
            \]
            Esto prueba que $f_n \xrightarrow{w^*} 0$.
            
        \item[Opción 2.]Sea $h\in L^1(\mathbb{R})$. Entonces
            \[
            \left| \int_{-\infty}^{+\infty} f_n(x)h(x)\,dx \right|
            = \left| \int_{-\infty}^{+\infty} \Phi(n+x)h(x)\,dx \right| = ({1})
            \]
            Como el soporte de $\Phi$ está contenido en $(0,1)$,
            \[
            \Phi(n+x)\neq 0 \Longrightarrow n+x\in(0,1)
            \Longrightarrow x\in(-n,-n+1).
            \]
            Por tanto,
            \[
            (1) = \left| \int_{-n}^{-n+1} \Phi(n+x)h(x)\,dx \right|
            \le \|\Phi\|_{L^\infty(\mathbb{R})}
            \int_{-n}^{-n+1} |h(x)|\,dx.
            \]
            Es decir, notando $I$ a la funcion indicadora
            \[
            (1) \le \|\Phi\|_{L^\infty(\mathbb{R})}
            \int_{-\infty}^{+\infty}
            I_{(-n,-n+1)}(x)|h(x)|\,dx. 
            \]
            
            Observamos que
            \[
            I_{(-n,-n+1)}(x)|h(x)| \xrightarrow[n \to \infty]{}0,
            \qquad \forall x\in\mathbb{R},
            \]
            y además
            \[
            |I_{(-n,-n+1)}(x)h(x)| \le |h(x)| \in L^1(\mathbb{R}).
            \]
            Por el teorema de la convergencia dominada,
            \[
            \int_{-\infty}^{+\infty}
            I_{(-n,-n+1)}(x)h(x)\,dx
            \xrightarrow[n \to \infty]{}0.
            \]
            Luego
            \[
            \int_{-\infty}^{+\infty} f_n(x)h(x)\,dx
            \xrightarrow[n \to \infty]{}0,
            \]
            lo que prueba que $f_n \stackrel{\ast}{\rightharpoonup} 0$ en $L^\infty(\mathbb{R})$.
        \end{description}
    \end{enumerate}
\end{ejercicio}
\end{document}
